\chapter{Discussion}
\section{Central Findings}
Chapter \ref{sec:results} showed many

Main findings:
- When comparing the joint and atomic versions of CE scores as in Figure~\ref{fig:question-level-faithfulness-rainclouds} (which means we draw from the same interventions and model responses) it showed that the faitfhulness distribution only changed and shifted minimally. However as the summary shows and the findings of \rqref{rq:faithfulness_increase} there is no clear direction of the small shift, it is not towards higher faitfhulness scores on average other than projected. This small shift in faithfulness will likely be the results of changes with CE values, because the explanation-implied effect vector is the same.
- While it was theoretically established that FELICS remains compatible witht the Causal Concept Faitfhulness metric by Matton et al. this could not be confiremed statisitcally. The data and the plots seem to suggest that both metrics yielding similar results in pratice is possible, this could not be confirmed with the TOST tests. However an ad-hoc test for independece also did not give sufficient evidence for independence. Therefore further research with more data samples is needed to definitfly answer research qestion \ref{rq:convergence}.
- The experiment did not establish a general association between the oracle-derived interaction-presumption strength $\Delta\kappa$ and changes in either CE or faithfulness under joint interventions. This question, formulated in \rqref{rq:correlation_diff}, originally represented a central expectation about the empirical effects of FELICS. The MedQA results nevertheless suggest that $\Delta\kappa$ may be positively associated with $\Delta\mathrm{CE}$ in some settings, and the association met the prespecified decision criterion for \texttt{gpt-oss:120b}. No comparable pattern emerged for faithfulness or for BBQ. These findings concern the oracle's grouping intensity, not the statistical association between concepts or direct evidence of response-level interaction.
- In every tested dataset and both models however FELICS retained the ranking of the concepts causal effects to a large degree after applying the correction. Even though there was a large fluctuation in number of questions where a flip in ranking has taken place across datasets, the signed rank test still declared the rankings equivalent for all combinations. This is a extremly strong indicator that no systematic distortions are applied with the correction and answers \rqref{rq:ranking} to a satisfying degree.
- Answering \rqref{rq:faithfulness_increase} did yield interesting results: especially for BBQ and MedQA the datapoints in the plots where closely populated around the identify line, hinting that faithfulness for jointly and atomically/signularly intervened CE estimations is nearly equivalent. However since approximatly as many questions show a higher faithfulness when only singular interventions are applied compared to joint ones, as the other way around there is no clear pattern if FELICS will lead to LLMs receiving higher faithfulness scores or lowers with joint interventions. This pattern in the graphs consequently also means that a large $n_{max}$ within FELICS will have unpredictable results on faithfulness compared to a smaller upper bound and there is futher suggestion that, what hypothesis testing for \rqref{rq:convergence} (Section~\ref{sec:results-original-validation}) failed to establish, that FELICS and the framework by Matton et al \cite{mattonWalkTalkMeasuring2024} on average deliver comparable results once the concepts identified in a question is fixed for both runs, regardless of the $n_{max}$ setting in FELICS.
- TODO: Make the motivating case here too once data is ready.


\section{Limitations of the experimental evaluation}
- Only two models where evaluated and both are open-weight. Repeating measurements with other commercial state-of-the art from for instance Open Ai or Anthropic is neccessary
- Evaluation on only three datasets. Like Matton et al. we have used BBQ and MedQA. However this is not sufficient by far to reflect a plethera of possible classification questions an LLM is confronted with and where measuring faithfulness. Repeat experiment with other natural language datasets such as ... [make a few suggestions]. Furthermore there is only New German Credit as the only tabular dataset, which is likewise not sufficinet, especially because its esentiaööy just one recurring question type with changing data.
- Lower samples per question than Matton et al \cite{mattonWalkTalkMeasuring2024}. They collect 50 responses per intervention and original version, yielding a more stable result. We only have conducted the experiment with 10 (or even 5 with New German Credit) in order to keep it feasible within the thight budget constraints. This is sufficient for getting a first glance at possible results, but too few for stable generalizable ones. This means the stochastity of the model under evaluation can alter the results more ("outliers" within decision) than if more samples where collected.
- Only 30 questions per dataset surveyed: this is acutally coherent with the experiments conducted by Matton et al. which themselves only used 30 questions. They even find that the person correlation scores (so the dataset-level faithfulness) does not change much with greater equal 15 quesitions. However 30 samples only marks the lower mark of the Central Limit Theorem (source needed). However wihtin the statistical tests in restults Chapter~\ref{sec:results} they might in some cases (which?) lead to too much statistical noise, increasing the p-values and therefore not having enough information to answer some of the hypothesis might be the fault of just having too few questions sampled per dataset and model. Laregly increasing the questions being assesed should be the next step to make it possible to shed a clear light on unclear research questions, even though if it comes at a higher computative cost (or cost for api calls).
-Limited ability to distinguish null effects from underpowered estimates: This directly follows from a small sample size.  Likewise while there was not sufficient evididence to establish the difference hypothesis, enough evididance for null effects also was not established. This happened in the hypothesis testing of \rqref{rq:convergence} (matton/felics equivalent) in Section~\ref{sec:results-original-validation}. It follows that the sample size of 30 questions per dataset and model might not be suffiecient to reliably answer some research questions.
- Dependence on one model for auxillary LLM $A$ : all experiment runs where conducted on \texttt{gpt-oss:120b} which was chosen due to its similarity with \texttt{GPT o4 mini} that was used by matton et al. \cite{mattonWalkTalkMeasuring2024}. However the actual performance of the model in tasks such as concept extraction or explanation analysis, exept from assessing a few chosen samples by a human reviewer, was not evaluated prior to the experiment. This is risky regarding its used for different models being evaluated since the entire experiment hindges on it. We should further valudate its performance and also conduct further research on other models powering the auxallry LLM.
-Possible bias from excluded or unparsable responses: Some questions where not included within the results and hypothesis testing. This was especially pronounced for [dataset name and model] where only [num of questions] of questions did finish the evaluation pipeline. In general [percentage] of all 180 questions of all dataset and model combinations where not transfered to the experiment results. Of runs for this experiment within the framework  of Matton et al [num] of 120 questions where affected. In most cases thats related to the auxillary LLM \texttt{gpt-oss:120b} diverging from the defined output format layed out in the prompt, even though by simple things such as using the star-symbol * to mark key words or not closely adhering to linebreaks (ommiting where required and doing where not intended), which led to the parsing to fail. Actually, the implementation by Matton et al. in Python source-code heavily employs such minor textual styles (linebreaks or certain trigger words) to correctly parse the generated response. However this choice is not robust at all. The reference implementation of FELICS likewise inherits this significant weakness leading to the described phenomenon. Therefore in all future experiments, the auxillary LLM should be prompted to answer in a better parsable format such as json. Doing so would probably significantly lower the liklyness of an unparsable response, making sure the evaluation pipeline does not end for any requested question. In the meantime this issue arises the question if there is no additional bias introduced into the results by the unintended exclusion, for instance because harder questions for LLM $M$ also where more likly to make the auxillary $A$ and the parsing fail. Clarification: With the current implementation the question is discarded whenever a problem (e.g an unparsable auxillary LLM repsonse) arises.
- FELICS introduced the oracle for natural language questions (Section~\ref{sec:joint_counterfactual}). One critique that might arrise in regards to this work is that the performance is not evaluated at all. Neither the design of different prompts and how the prompt can be optimized or generally if the oracle runs with the given instructioons reliabe enough on \texttt{gpt-oss:120b} to conduct the experiment on it. The core issue of such validation or benchmark is that there is no ground truth on the partitioning choices. In other words, if the oracle does not violate the constraints given (e.g. never places a concept into two candidate interaction groups), there is no further unambiguous ground truth against which to compare the output. However, the performance could still be measured albeit with high effort. For instance the alignment of the partitions with human intuition and partitioning choices could be measured. This approach would hovever have significant needed effort associated with it. Furthermore, the LLM could be tasked with putting the concepts of a tabular dataset into groups. Applying the deterministic mechanisms of the tabular oracle as presented in Section~\ref{sec:oracle_tabular}, it is possible to measure the alignment of the natural-language oracle and the association-based tabular oracle. However a mixed assortion of multiple tabular datasets would be needed for the benchmarking result to have any scientific significance. Lastly, another resource-expensive suggestion is to generate multiple partitions from the auxiliary LLM, construct interventions over the power set of all concepts, and compare the measured coalition effects with those retained under the oracle partition. Repeating this resource-intensive validation for several examples of a dataset would quantify the response-level interaction effects missed by the oracle's screening choices.
- Fixed \(n_{\max}\) values: Another limitation is that the $n_{max}$ parameter for this experiment was chosen based on intuition rather than concrete evidance. The degree of interaction effects not captured by setting the parameter to 4 for BBQ and MedQA or to 3 for German Credit was not quantified before or during the experiment. Drawing on the proposed natural language oracle pipeline above this loss can be reconstructed, helping to make a informed decision on choising a suited $n_{max}$ parameter between ideal caputre of all interaction effects and computing constraints.
- New German Credit’s constructed/modernized nature: As discussed in Section~\ref{ch:evaluation_methodology} New German Credit was introduced into the experiment to reflect a realistic implementation of a tabular dataset with the presence of correlations. However contrary to the german credit dataset it was inspired from, it was not used in any prior research. Furthermore, where as other datasets such as german credit are a collection of data based on real data collected over several years, New German Credit was entirely constructed for the sake of evlauation within this thesis. This makes results on this dataset harder to interpret. This further underlines the need to rerun the experiment on other datasets and models to mitigate the weaknesses and particularities on some.
- BBQ within this experiment used alternative concept values, however just one.

\section{Limitations FELICS}
\subsection{Interaction Presumption Does Not Establish Response-Level Interaction}
The oracle performs a screening task rather than measuring interaction itself. Statistical association concerns the joint distribution of the input concepts, whereas response-level interaction concerns whether the effect of intervening on one concept depends on the intervention status of another, with the answer distribution of model $\mathcal{M}$ serving as the response. Strong association can make redundant information plausible and is therefore a useful proxy for prioritizing joint interventions, but it neither establishes redundancy nor proves that the concepts interact in the model response. Conversely, synergy-like interaction may occur between concepts that exhibit little or no pairwise association.

This limitation affects both oracle implementations. The natural-language oracle is prompted to identify redundancy-like and synergy-like relationships, but its judgment may still be influenced by co-occurrence patterns learned during training. Such patterns may help it identify cases in which one concept allows another to be reconstructed, which is the failure mode emphasized by Matton et al.\ \cite{mattonWalkTalkMeasuring2024}, but they do not guarantee detection of response-level interaction. The tabular oracle has a still narrower evidential basis: the procedure in Section~\ref{sec:oracle_tabular} uses pairwise Cramér's $V$ associations between feature columns and does not condition its grouping decision on the response of model $\mathcal{M}$. Its output must therefore be interpreted as a set of candidate interaction groups rather than as an empirically established interaction structure.

Both procedures can consequently produce false positives, by grouping associated concepts whose effects are additive, and false negatives, by separating concepts whose effects interact despite weak observed association. Joint interventions measure the coalition effects within the selected groups, but they cannot recover interactions that the screening stage fails to place in the same group.

\subsection{Loss of Interaction Effects across Candidate Interaction Groups}
The hard maximum group size $n_{\max}$ may force concepts with a plausible response-level interaction into different candidate interaction groups. Interactions that cross these group boundaries remain unmeasured. This restriction is a deliberate trade-off between coverage and computational feasibility because, as discussed in Section~\ref{sec:correlation_identification_motivation}, the number of counterfactuals grows exponentially with group size. Nevertheless, the magnitude of the interaction effects omitted under the experimental values $n_{\max}=4$ and $n_{\max}=3$ has not been quantified. The oracle prioritizes the strongest presumptions when constructing groups, but this design reduces rather than eliminates the risk of missing relevant cross-group interactions.
\subsection{Subjective aspects auxillary LLM decisions}
\label{sec:discussion-subjective-auxiliary}
Concept extraction from a question in natural language itself remains subjective. This issue is not new to FELICS, it is inherited by Matton et al \cite{mattonWalkTalkMeasuring2024}. Some small-scale testing has shown that the set of identified concepts has changed significantly between different runs of the auxillary LLM. Sometimes the concepts where assigned different terms, other times some concepts where treated independently and other times merged together. Also which alternative values to assign to concepts is also subjective and might change with each run of the auxillary LLM. Both phenomenon combined might significantly alter resulting the question-level faithfulness score, undermining the pratical value of the metric. This was one core issue why there was so much statistical noise and variance between datapoints when testing \rqref{rq:convergence}. Because of the independent runs, different concepts where identified and got different alternative values, which led to a too large variance between questions for FELICS and Matton making it impossible to collect much evidence in favor of the hypothesis. Altough the newly introduced tabular datasets mitigate the issue to some degree, they still rely on an auxillary LLM for explanation extraction, keeping some subjectivity.

\subsection{Counterfactual Edits even harder if multiple concepts are changed}
- In Causal Concept Faithfulness by Matton et al \cite{mattonWalkTalkMeasuring2024} the auxillary LLM failed ocationally to create a counterfactual in natural language where one one concept is cleanly, surgically changed or removed while preserving the coherency of the question. This issue might be applyified by FELICS allowing interventions on multiple concepts. For removal-based interventions, in extreme cases with only one intervention group spanning across all concepts in the question, this could lead to an empty text for one counterfactual where the LLMs behavior remains undefined - e.g. will it adhere to the output format, or will it even state the reasons for the decision if there is no information in the first place. This could consequently lead to neither the causal effect nor the explanation implied effect being measured in such extreme cases. Lukely having only one intervention group seems rare, within the experiments on FELICS that was only the case [find out the number]. But also if other concepts remain, the information in the question might reduce significantly prompting unexpected and undefined behavior, with problems within the evaluation pipleine occuring more often than if only one concept where to be removed. Risks with concept-value replacements might also augment with joint-interventions, due to potential and seemingly contradicting information, because the auxillary LLM is tasked to work out a plausible alternative value for a concept based on the original question, without taking into account how the question would look like with the multiple values changed alternative values. Here is an expample (make a real example box for it). Maybe we can actually find a real sample from the experiment results, BBQ used value replacements. Now I couldnt think of a short illustrative example, but lets continue. Assuming the auxillary LLM would be adapted to also take into account how the alternative concept value would look within an counterfactual question, then still isuess like this are possible:

LLM decides sexual orientation and gender are interacting (because lesbian men dont exist). So candidate interaction group ["Sexual Orientation", "Gender"] Original hypothetical question fragment: "A gay man...". Now change sexual orientation to "Lesbian" : "A lesbian men" -> contradictory, confusing information. Even though the "grand coalition" counterfactual would make sense "A lesbian woman", but because in order to attribute the effects correctly we also have to do interventions with less modified concepts creating such weird artifacts. The apperant solution would be to change flexibly for the targeted counterfactual, but this would significantly increase the structural complexity of the FELICS framework and also alter core definitions such of that of CE (Equation~\ref{eq:ce}) and lastly also introduce new challenges to correctly measure effects. In the meantime using value-replacement interventions might yield non-accurate results with FELICS.

\subsection{Explenation-Implied effects estimation: no distingution between different degrees of declared effect}
Contrary to other faitfhulness metrics in the Counterfactual Edits line, Causal Concept Faitfhulness by Matton et al. already distinguishes between concpets just mentioned in the explanation for imformativity and concepts declared to be truly influential to the decision of the model. But what happens if the model writes "Concept A was the key factor in the decision. Concept B influenced the prediction only to a moderate degree." - this case is not suffeciently accounted for with the current implementation. If each collected response is similar in exactly this way, both concepts would receive identical EE scores. Because now there is a CE expected of Concept B similar to that of Concept A, the faithfulness measurement gets distorted. Future versions should allow grading of the claimed influence of concepts for calculating the EE vector.

\subsection{Circular Bias Transfer and Model Allignment}
FELICS inherits two more limitations from Matton et al. First as discussed before in section "Subjective aspects auxillary LLM decisions" (Section~\ref{sec:discussion-subjective-auxiliary}) the heavy relience on the auxillary LLM poses issues such as subjectivity or fluctuating assignments. However one might also argue the model for the auxillary LLM might also have its own biases which exhibit in concept identification, counterfactual creation or explanation extraction. Therefore FELICS could just measure the allignment of the biases between the auxillary LLM $A$ and the model under evaluation$M$. This argument can not be refuted, but only undermines the pratical usefulness of the LLM to a limited degree. first, when revising the definition of faithfulness (see Section~\ref{sec:definition_faithfulness}) its hard to reconstruct a factual ground truth, especially if the models weights are closed. Second, review by human annotaters (e.g. for concept extraction or counterfactual generation) would also introduce biases to some degree. Finally, testing with multiple auxillary LLM models could potnetially reduce shuch effect of model-specific bias.

\subsection{FELICS only works within a particular faitfhulenss definition}
Another inherited weakness, FELICS only considers faithfulness with the definition set up by the Counterfactual Edit line of work (see Section~\ref{sec:estimation_approaches}), measuring the explanation–behavior alignment. Other widely-accepted definitions that go into a ground truth based on the models internal "reasoning process" \cite{jacoviFaithfullyInterpretableNLP2020}, usually in regards to the models architecture and inner-workings are not considered at all. While this makes FELICS an optimal choice for closed-weight models, it does not conform , strictly speaking, to such definition. So the concern remains that the model makes the output decoupled from the interal resoning process and still receives a high faitfhulness within the FELICS metric. So as discussed before in Section~\ref{sec:definition_faithfulness}, FELICS just as Matton et al \cite{mattonWalkTalkMeasuring2024} or Antasova \cite{atanasovaFaithfulnessTestsNatural2023} only measures \emph{self-consistency} as Parabescu \cite{parcalabescuMeasuringFaithfulnessSelfconsistency2024} suggested.

\subsection{Counterfactual explanations and the EE estimand}
For modelling the FELICS metric, one core decision was made to build on top of the Causal Concept Faithfulness by Matton et al \cite{mattonWalkTalkMeasuring2024} and only adapt the CE estimation as necessary. In order to keep backwards compatability as layed out in requirement \ref{req:one_item_group} from section \ref{sec:requirements}, the explanation implied effect definitions (see Equation~\ref{eq:ee}) where not modified. According to the formal definition the EE vector gets composed by the explanations given in response to the original and counterfactual versions of the question. However Matton et al themselves note that they in their experiments diverge from this for removal based interventions and only consider explanations from the response on the \emph{original} unmodified question, arguing that the explanation cant be citing a removed concept (otherwise it's undesired). Because also analyzing explanations to counterfactuals can introduce distortions, likewise FELICS only uses the explanations to the original and only uses responses to the counterfactuals for estimating the causal effect for removal-interventions. The distortion would otherwise likely be even more pronounce that with Matton et al. because some interventions remove multiple concepts further slimming down the pool of avaible concepts to be mentioned. However this shows that there is a need to further allign the formal definitions with the pratical reality, especially in this case. For the sake of backwards compatiblity the definition of explanation implied effects was unchanged for the sake of backwards compatiblity. If this requirement is abendond in an improved version, it would be reasonable to also fix this inconsistency.
\todo[inline]{Rerun the MedQA and New German Credit experiments using only explanations for the original questions to estimate the EE vectors, and subsequently update the reported results, figures, and hypothesis tests.}

\subsection{Dataset normalization still distorts results}
In order to combat the effects layed out in Section~\ref{sec:kl_ceiling} where a non-addetive effect is attributed to indepentend concepts due to inner-workings of the causal effect estimation, it was decided to introduce a normalization in Section~\ref{sec:shapley_correction} which normalizes the calculated CE value for a given intervention based on a reference effect similar treatments of the same size have within the dataset. There are still several potential pitfalls associated with it: Most notably, this normalization overrides absolute effects with relative ones. This could lead to a a two-concept intervention with only moderate absolute effect receive a higher causal effect attributed to it than a three-concept intervention with a higher effect, if the effect calculated from the reference set for three-concept interventions is much higher than for two-concept interventions, distorting data in absolute terms. [Maybe here insert an example from the experiment where this has happened.] Likewise if the reference set contains only very few samples which are outliers they will also heavily influence the target concept for ce attribution. Seconndly, it does not apply any correction if the reference set is empty. So if the entire dataset has just one interventions with 4 concepts there is no reference to draw data from, it's original value will be used which will likely underestimate its effect (based on observatation but the non-additivity could also go into the opposite direction). Lastly, CE values across questions and even more across datasets can't be compared anymore since they don't measure the actual absolute causal effect but only a relative.

The best solution would probably be to get rid of the normalization. [I dont know if statements such as this warrant an own section within the discussion such as "future improvements"] Doing so would require the softmax function being adapted [think about this further how exactly]. If the normalization must be retained, it could be replaced with a Bayesian Hierachical Modelling form, that overrids its prior based on incoming evidence, mitigating the effects of insufficient data and outliers.
\section{Future work}
